\section{Methods}
The prediction unit is the participant because the target label is participant-level.
Word-level rows provide repeated observations for estimating gaze profiles, but they
are not independent prediction targets. The primary validation split is
leave-one-participant-out (LOPO): each fold trains on 56 participants and predicts the
single held-out participant, producing 57 predictions and zero skipped folds.

DFM base-model features are extracted at the word level using surprisal and entropy as
contextual predictability measures. The locked primary model is
\texttt{D3\_dfm\_residual\_gaze\_only}: a standardized logistic regression over
participant-level residual gaze sensitivity features. No random word-level split is
used, and no deep model or broad feature expansion is introduced.

Cross-fitted residualization is performed inside the held-out participant split
(Figure~\ref{fig:crossfit-schematic}). For each LOPO fold, expected gaze models are fit
using only the training participants. Predictors include lexical, positional, DFM, text,
boundary-opacity, and quality covariates available in the frozen package. Those models
predict expected gaze for the held-out participant's word rows, and held-out residuals
are aggregated into participant-level DFM surprisal and entropy sensitivity slopes.
Reader group is never used in residualization, and the held-out participant's rows are
never used to fit that fold's residual model.

The confirmatory comparisons separate DFM exposure from DFM sensitivity. Exposure-only
features summarize the text a participant saw; sensitivity features summarize how a
participant's residual gaze costs vary with DFM predictability. The primary model also
excludes direct exposure-count variables, including number of words read, number of
speeches, number of word rows, total word rows, and word-observation count. Ablations
remove exposure-count variables and raw-speed/global-duration features. Robustness is
evaluated with label permutation, participant bootstrap intervals, calibration
statistics, influence analysis, and coefficient sign stability. Focused interaction
models are used only for interpretation of reader-group differences involving word
length, DFM surprisal, and previous-boundary opacity.
